% Contribution to the polylogarithms continuation.
% Baseline: cc34f73596336f2466d9754cb0f3635bd2bedade.
% Parent document supplies amsmath, amssymb, amsthm, theorem environments.
\section{A uniform transition for two large reflected log-gamma exponents}
\label{joint:sec:moments}

Write
\[
 f(x)=\log\Gamma(x),\qquad
 M_{n,m}=\int_0^1 f(x)^n f(1-x)^m\,dx,
 \qquad n,m\in\mathbb Z_{\geq0}.
\]
The mixed integrals themselves are classical: they occur as
$T_{n+m,n}$ in Amdeberhan et al.\ \cite[Section~5]{gaussian:ACEMKM}.
The pinned manuscript proves a fixed-$m$ expansion, its late coefficients,
and a sharp remainder at optimal truncation. It explicitly leaves the
regime in which both exponents increase as a separate question.
Here we resolve a transition inside that joint regime.

For fixed $m$, the first endpoint term is
\[
 M_{n,m}\sim\frac{\gamma^m n!}{(m+1)^{n+1}}.
\]
When $m$ increases, its first neglected local perturbation has size
$m\exp[-n/(m+1)]$. Thus it can survive in the limit even though the
relevant endpoint $x$ tends to zero. The transition occurs around
$n/m=\log m$. The theorem below gives its complete expansion, with
explicit first and second corrections. In particular, it does not
substitute a growing $m$ into a fixed-$m$ error estimate.

\subsection{The transition theorem and its first two corrections}

Define the exact constants
\begin{align}
 c&=\frac{\zeta(2)}{2\gamma},&
 d&=\frac{\zeta(3)}{3\gamma}-\frac{c^2}{2},&
 e&=\frac{\zeta(4)}{4\gamma}
       -\frac{c\zeta(3)}{3\gamma}+\frac{c^3}{3},\label{joint:eq:constants1}\\
 \delta&=c-\gamma,&
 \varepsilon&=\gamma^2+d,&
 \eta&=e+\frac{\gamma^3}{2}-\frac{\gamma^2\delta}{2}
          +\frac{\gamma\delta^2}{2}-5\gamma\varepsilon.
 \label{joint:eq:constants2}
\end{align}
In particular,
\[
 \delta=0.84767122475766887678\ldots,
 \qquad
 \varepsilon=0.01219631601207367653\ldots.
\]
The definitions, rather than these diagnostic decimals, specify the
constants in all statements. Notice that $\delta>0$; for example,
$\gamma<0.6$ and $\zeta(2)>1.6$ already suffice.

\begin{theorem}[Uniform joint-exponent transition]\label{joint:thm:transition}
Let
\[
 r=m+1,\qquad t=\frac{n+1}{m+1},\qquad
 \lambda=me^{-t},\qquad
 \mathcal R_{n,m}=\frac{(m+1)^{n+1}M_{n,m}}{\gamma^m n!}.
\]
Fix $0<a<b<\infty$. As $m\to\infty$, uniformly over positive
integers $n$ for which $a\leq\lambda\leq b$,
\begin{equation}\label{joint:eq:two-corrections}
 \mathcal R_{n,m}
 =e^{\delta\lambda}
 \left\{1+\frac{P_1(t,\lambda)}m
          +\frac{P_2(t,\lambda)}{m^2}
          +O_{a,b}\left(\frac{t^3}{m^3}\right)\right\},
\end{equation}
where
\begin{equation}\label{joint:eq:P1}
 P_1(t,\lambda)=\frac{t}{2}\delta\lambda(1+\delta\lambda)
                      -2\gamma\lambda+\varepsilon\lambda^2,
\end{equation}
and
\begin{align}
 P_2(t,\lambda)={}&
 \lambda\left\{\frac{\delta t^2}{8}
                    -\left(\gamma+\frac{5\delta}{6}\right)t\right\}
 \notag\\
 &+\lambda^2\left\{\frac{7\delta^2t^2}{8}
        +\left(-3\gamma\delta-\frac{3\delta^2}{2}
                        +2\varepsilon\right)t
        +\frac{15\gamma^2}{2}+3\gamma\delta\right\}\notag\\
 &+\lambda^3\left\{\frac{3\delta^3t^2}{4}
        +\left(-\gamma\delta^2-\frac{\delta^3}{3}
                           +\frac{5\delta\varepsilon}{2}\right)t
        +\eta\right\}\notag\\
 &+\lambda^4\left\{\frac{\delta^4t^2}{8}
                    +\frac{\delta^2\varepsilon t}{2}
                    +\frac{\varepsilon^2}{2}\right\}.
 \label{joint:eq:P2}
\end{align}
More generally, for every fixed integer $J\geq0$, there are explicitly
computable polynomials $P_0=1,P_1,\ldots,P_J$ such that
\begin{equation}\label{joint:eq:all-orders}
 \mathcal R_{n,m}=e^{\delta\lambda}
 \left\{\sum_{\ell=0}^J\frac{P_\ell(t,\lambda)}{m^\ell}
       +O_{a,b,J}\left(\frac{t^{J+1}}{m^{J+1}}\right)\right\}.
\end{equation}
They have degree at most $\ell$ in $t$ and at most $2\ell$ in
$\lambda$, and
\[
 P_\ell\in
 \mathbb Q[\gamma,\gamma^{-1},\zeta(2),\ldots,\zeta(\ell+2)]
 [t,\lambda].
\]
The finite coefficient prescription in
\eqref{joint:eq:polynomial-prescription} below determines every
$P_\ell$.
\end{theorem}

Here $t=\log m-\log\lambda=\log m+O_{a,b}(1)$, so the displayed
remainders are successive powers of $(\log m)/m$.

\begin{corollary}[The surviving correction to the endpoint formula]
\label{joint:cor:limit}
If $m\to\infty$ and
\[
 \frac{n+1}{m+1}-\log m\longrightarrow s\in\mathbb R,
\]
then
\begin{equation}\label{joint:eq:limit}
 \frac{(m+1)^{n+1}M_{n,m}}{\gamma^m n!}
 \longrightarrow
 \exp\left[\left(\frac{\zeta(2)}{2\gamma}-\gamma\right)e^{-s}\right].
\end{equation}
In particular the fixed-$m$ leading term does not have uniform relative
error $o(1)$ on this window. At $s=0$ its missing factor tends to
$e^\delta=2.3342\ldots$.
\end{corollary}
\begin{proof}
Use $\lambda\to e^{-s}$ in \eqref{joint:eq:all-orders} with $J=0$.
The corrections tend to zero because $t/m\to0$.
\end{proof}

\subsection{Exact gamma averaging and localization at the correct endpoint}

Put
\[
 g(x)=\log\Gamma(1+x),\qquad
 B(x)=\log\Gamma(1-x),\qquad
 C(x)=\frac{B(x)}{\gamma x},\quad C(0)=1.
\]
The classical convergent log-gamma expansion
\cite[\S5.7]{joint:DLMF} gives
\begin{align}
 g(x)&=-\gamma x+\frac{\zeta(2)}2x^2
                         -\frac{\zeta(3)}3x^3+O(x^4),\notag\\
 \log C(x)&=cx+dx^2+ex^3+O(x^4).
 \label{joint:eq:local-series}
\end{align}
All these germs are holomorphic in a fixed disk about zero.
For $0<x<1$, log-convexity of $\Gamma$ between $1$ and $2$ implies
\begin{equation}\label{joint:eq:elementary-envelopes}
 0<f(x)\leq-\log x,\qquad
 0<B(x)\leq-\log(1-x).
\end{equation}
After $x=e^{-y}$, the normalized integral is the exact expectation
\begin{equation}\label{joint:eq:expectation}
 \mathcal R_{n,m}=\mathbb E e^{S(Y)},\qquad
 S(y)=n\log\left(1+\frac{g(e^{-y})}{y}\right)
                         +m\log C(e^{-y}),
\end{equation}
where $Y$ has a gamma distribution with shape $q=n+1=rt$ and rate
$r=m+1$. Thus
\[
 \mathbb EY=t,\qquad
 \operatorname{Var}Y=\frac tr.
\]
The powers and logarithms in \eqref{joint:eq:expectation} are real,
since the original factors are positive.

\begin{lemma}[Uniform localization]\label{joint:lem:localization}
Under the compact-$\lambda$ hypotheses of
Theorem~\ref{joint:thm:transition}, there is $\kappa>0$ such that
\begin{equation}\label{joint:eq:localization}
 \mathbb E\left[e^{S(Y)}\mathbf1_{|Y-t|>1}\right]
       =O_{a,b}\left(e^{-\kappa m/t}\right).
\end{equation}
The same estimate holds with $e^{S(Y)}$ replaced by any fixed power
of $|Y-t|$. In particular, the endpoint $x\to1$, corresponding to
$y\to0$, is negligible on the normalization used in the theorem.
\end{lemma}
\begin{proof}
The first summand of $S(y)$ is nonpositive by
\eqref{joint:eq:elementary-envelopes}. Analyticity of $C$ and $C(0)=1$
give a constant $K$ such that
\[
 e^{S(y)}\leq\exp(Kme^{-y})\qquad(y\geq1).
\]
For $y\geq t+1$ this weight is uniformly bounded, and the gamma
Chernoff bound gives a tail $O(e^{-\kappa m/t})$.
For example the gamma rate function is $I(u)=u-1-\log u$, and
$rt I(1+1/t)\geq r/(3t)$ for all sufficiently large $t$.

For the lower tail, let $p(y)$ denote the gamma density and put
$v=t-y$. When $1\leq y\leq t-1$, we have
\[
 \frac{p(t-v)}{p(t)}
  =\exp\{(rt-1)\log(1-v/t)+rv\}
  \leq t\exp\left(-\frac{rv^2}{2t}\right).
\]
Here $1\leq v\leq t-1$, and $me^{-(t-v)}=\lambda e^v$.
The maximum of $e^v/v^2$ on this interval is at one of its endpoints;
consequently
\[
 K\lambda e^v\leq\frac{rv^2}{4t}
 \qquad(1\leq v\leq t-1)
\]
for all sufficiently large $m$, uniformly in $a\leq\lambda\leq b$.
Indeed the two required bounds reduce to $O(1)\leq r/(4t)$ and
$O(m/t^2)\leq r/(4t)$. It follows that the weighted lower-tail
integral over $[1,t-1]$ is at most
\[
 p(t)t^2e^{-r/(4t)}.
\]
Stirling's inequalities give $p(t)=O(\sqrt{r/t})$; the polynomial
factor is absorbed by slightly reducing $\kappa$.

It remains to bound $0<y<1$. Since $1-e^{-y}\geq y/e$ there,
\[
 B(e^{-y})\leq1-\log y,
 \qquad (1-\log y)^m\leq y^{-m}.
\]
Using $f(e^{-y})\leq y$ in the original normalized integral bounds
this part by
\begin{equation}\label{joint:eq:opposite-endpoint}
 \frac{r^{n+1}}{\gamma^m n!}
 \int_0^1 y^n(1-\log y)^m\,dy
 \leq\frac{r^{n+1}}{\gamma^m n!(n-m+1)}.
\end{equation}
Since $n+1=rt$ and $t=\log m+O(1)$, the logarithm of this bound is
\[
 rt(1-\log t)+O(m+\log n),
\]
which is smaller than $-\kappa m/t$ for large $m$.
This also shows directly why the other endpoint cannot compete.
The unweighted gamma estimates, with polynomial factors, follow by
the same density estimates or by Chernoff's inequality with a
slightly smaller exponent.
\end{proof}

\subsection{Derivation of the first two corrections}

Write $U=Y-t$. On $|U|\leq1$, set $x=(\lambda/m)e^{-U}$ and use
\eqref{joint:eq:local-series}. Uniformly there, together with every
fixed number of $U$-derivatives,
\begin{equation}\label{joint:eq:S-expansion}
 S(t+U)=S_0(U)+\frac{S_1(U)}m+\frac{S_2(U)}{m^2}+O(m^{-3}),
\end{equation}
where
\begin{align}
 S_0(U)&=\lambda e^{-U}
       \left(c-\frac{\gamma t}{t+U}\right),\notag\\
 S_1(U)&=-\frac{\gamma\lambda e^{-U}(t-1)}{t+U}
  +\lambda^2e^{-2U}
       \left(d+\frac{\gamma ct}{t+U}
                   -\frac{\gamma^2t}{2(t+U)^2}\right),\notag\\
 S_2(U)&=\lambda^2e^{-2U}(t-1)
       \left(\frac{\gamma c}{t+U}
                    -\frac{\gamma^2}{2(t+U)^2}\right)\notag\\
 &\quad+\lambda^3e^{-3U}
       \left(e-\frac{\gamma t(d+c^2/2)}{t+U}
          +\frac{\gamma^2ct}{(t+U)^2}
          -\frac{\gamma^3t}{3(t+U)^3}\right).
 \label{joint:eq:S012}
\end{align}
The remainders are uniform because $t/(t+U)$ is bounded,
$x=O(m^{-1})$, and $n=O(mt)$. The displayed amplitudes and all their
fixed-order derivatives are uniformly bounded.

The central gamma moments needed for the first correction are
\[
 \mathbb EU=0,\qquad
 \mathbb EU^2=\frac tr,\qquad
 \mathbb EU^3=\frac{2t}{r^2},\qquad
 \mathbb EU^4=\frac{3t^2}{r^2}+\frac{6t}{r^3}.
\]
In particular, for a smooth amplitude $H$ whose first four derivatives
are uniformly bounded on $[-1,1]$, Taylor expansion through the cubic
term, followed by the localization lemma, gives
\[
 \mathbb EH(U)=H(0)+\frac{t}{2m}H''(0)
                        +O(t^2/m^2).
\]
This notation means that $H$ is averaged on the localized interval;
any uniformly bounded extension outside is immaterial up to the
exponentially small error already bounded. The signed third moment is used: an absolute
third-moment estimate would only give an error of order
$(t/m)^{3/2}$.

Since $S_0(0)=\delta\lambda$, the first normalized correction equals
\[
 S_1(0)+\frac t2\bigl(S_0''(0)+S_0'(0)^2\bigr).
\]
The derivatives are
\[
 S_0'(0)=\lambda\left(-\delta+\frac\gamma t\right),\qquad
 S_0''(0)=\lambda\left(\delta-\frac{2\gamma}t
                                   -\frac{2\gamma}{t^2}\right).
\]
Substitution gives \eqref{joint:eq:P1}. The terms involving $1/t$
cancel exactly; this is an algebraic cancellation, not an additional
large-$t$ approximation.

For an independently reproducible second-order calculation, define
\[
 a_j=S_0^{(j)}(0)\quad(1\leq j\leq4),\qquad
 b_j=S_1^{(j)}(0)\quad(0\leq j\leq2).
\]
Then the second correction is exactly
\begin{align}
 P_2={}&S_2(0)+\frac{b_0^2}{2}
       -\frac t2(a_2+a_1^2)\notag\\
 &+\frac t2\{b_2+2a_1b_1+(a_2+a_1^2)b_0\}\notag\\
 &+\frac t3(a_3+3a_1a_2+a_1^3)\notag\\
 &+\frac{t^2}{8}
     (a_4+4a_1a_3+3a_2^2+6a_1^2a_2+a_1^4).
 \label{joint:eq:P2-derivative}
\end{align}
Expanding this finite formula gives \eqref{joint:eq:P2}.
To justify the remainder, Taylor-expand the leading amplitude through
degree five and bound its degree-six remainder. The central moments
satisfy
\[
 \mathbb EU^5=\frac{20t^2}{r^3}+\frac{24t}{r^4},\qquad
 \mathbb EU^6=\frac{15t^3}{r^3}+\frac{130t^2}{r^4}
                                  +\frac{120t}{r^5}.
\]
Use also
$\mathbb EU^2=t/m-t/m^2+O(t/m^3)$ and
$\mathbb EU^3=2t/m^2+O(t/m^3)$.
The $m^{-1}$ amplitude needs Taylor terms only through degree three,
and the $m^{-2}$ amplitude only through degree one with a quadratic
remainder. All omitted terms are $O(t^3/m^3)$.
This proves \eqref{joint:eq:two-corrections} with its claimed uniformity.

\subsection{Exact offset-residue identities and a finite all-order recipe}

For comparison with the fixed-$m$ manuscript, retain its coefficients
\[
 A_{k,m}=\frac1k[x^{k-1}]\Gamma(1+x)^kB(x)^m.
\]
Set
\[
 H(x)=g(x)+\log C(x)=\sum_{j\geq1}h_jx^j,
 \qquad g(x)=\sum_{j\geq1}g_jx^j.
\]
Then
\[
 h_1=\delta,\quad h_2=d+\gamma c,\quad
 h_3=e-\gamma(d+c^2/2),\qquad g_1=-\gamma.
\]

\begin{proposition}[Exact polynomial structure at every fixed offset]
\label{joint:prop:offset}
For $j,m\geq0$,
\begin{equation}\label{joint:eq:offset-identity}
 Q_j(m):=\frac{(m+j+1)A_{m+j+1,m}}{\gamma^m}
       =[x^j]\exp\{mH(x)+(j+1)g(x)\}.
\end{equation}
In particular $Q_j$ is a polynomial in $m$ of degree $j$, with
leading coefficient $\delta^j/j!$. Its first two terms, for $j\geq2$,
are
\begin{equation}\label{joint:eq:offset-leading}
 Q_j(m)=\frac{\delta^j}{j!}m^j
 +\left\{\frac{h_2\delta^{j-2}}{(j-2)!}
       -\frac{\gamma(j+1)\delta^{j-1}}{(j-1)!}\right\}m^{j-1}
 +O_j(m^{j-2}).
\end{equation}
For $j=0,1$ the identity gives respectively
$Q_0=1$ and $Q_1=m\delta-2\gamma$.
\end{proposition}
\begin{proof}
Write $B(x)^m=(\gamma x)^mC(x)^m$ in the coefficient formula
and use $m+j+1=m+(j+1)$. This gives
\eqref{joint:eq:offset-identity} exactly.
Each factor of $m$ in the exponential consumes at least one degree
of $x$, proving the degree assertion. The degree-$j$ term uses only
$h_1x$. The degree-$(j-1)$ term either uses $h_2x^2$ once or
$(j+1)g_1x$ once. These are precisely
\eqref{joint:eq:offset-leading}.
\end{proof}

The same identity provides a finite prescription for every correction
in Theorem~\ref{joint:thm:transition}. Let $z,w$ be formal variables and
put, for $j\geq0$,
\begin{align*}
 \Theta_j(z,t)
 &=t\sum_{r\geq1}\frac{(-1)^{r+1}j^{r+1}}{r+1}
                                  \left(\frac z{1+z}\right)^r,\\
 K_j(z,w)
 &=\exp\left\{\delta w+
      \sum_{r\geq1}z^r
        \bigl(h_{r+1}w^{r+1}+(j+1)g_rw^r\bigr)\right\}.
\end{align*}
Then
\begin{equation}\label{joint:eq:polynomial-prescription}
 \boxed{\displaystyle
 P_\ell(t,\lambda)=e^{-\delta\lambda}[z^\ell]
   \sum_{j\geq0}\lambda^j e^{\Theta_j(z,t)}[w^j]K_j(z,w).}
\end{equation}
Coefficient extraction in $z$ is performed before summing the
resulting entire series in $\lambda$. This is essential: no convergence
of the original infinite pole expansion is asserted.

\paragraph{Why the prescription is finite.}
At order $z^\ell$, only $h_1,\ldots,h_{\ell+1}$ and
$g_1,\ldots,g_\ell$ can occur. Every summand is a finite linear
combination of
\[
 j^p\frac{\delta^{j-d}}{(j-d)!}\lambda^j,
\]
with the expression understood to be zero for $j<d$.
Their sums are obtained by applying powers of
$D_\lambda=\lambda\,d/d\lambda$ to
$\lambda^d e^{\delta\lambda}$.
A term of cost $z^r$ in either exponential has total $j$-degree plus
$w$-degree at most $2r$. Therefore
$\deg_\lambda P_\ell\leq2\ell$; also
$\deg_tP_\ell\leq\ell$. The coefficients lie in the ring stated in
the theorem, since $h_{\ell+1}$ uses zeta values only through
$\zeta(\ell+2)$.

\begin{corollary}[Two coefficient identities at every asymptotic order]
\label{joint:cor:top-coefficients}
Let $T_k(u)=e^{-u}(u\,d/du)^ke^u$ be the Touchard polynomial,
equivalently $T_k(u)=\sum_{j=0}^k\genfrac{\{}{\}}{0pt}{}{k}{j}u^j$. Then
\begin{align}
 [t^\ell]P_\ell(t,\lambda)
   &=\frac{T_{2\ell}(\delta\lambda)}{2^\ell\ell!},\label{joint:eq:top-t}\\
 [\lambda^{2\ell}]P_\ell(t,\lambda)
   &=\frac1{\ell!}
           \left(\varepsilon+\frac{\delta^2t}{2}\right)^\ell.
 \label{joint:eq:top-lambda}
\end{align}
These are exact finite polynomial identities, for every $\ell\geq0$.
\end{corollary}
\begin{proof}
The highest possible power $t^\ell z^\ell$ in
\eqref{joint:eq:polynomial-prescription} comes only from the $\ell$th
power of $tzj^2/2$. Summing $j^{2\ell}\lambda^j\delta^j/j!$
gives \eqref{joint:eq:top-t}.
For the second assertion, every exponent term of order $z^r$ with
$r\geq2$ has total $j$-degree plus $w$-degree at most $r+1<2r$.
Therefore the degree $2\ell$ in $\lambda$ comes only from the
$\ell$th power of
\[
 h_2w^2+(j+1)g_1w+\frac t2j^2.
\]
At the highest degree, the summation rule replaces a monomial
$w^dj^p$ by $\delta^p\lambda^{d+p}$.
The resulting coefficient is
$(h_2+g_1\delta+t\delta^2/2)^\ell/\ell!$.
Since $h_2+g_1\delta=\varepsilon$, this is
\eqref{joint:eq:top-lambda}.
\end{proof}

\paragraph{Why these are the analytic asymptotic coefficients.}
The localization proof and \eqref{joint:eq:S-expansion} extend to any
fixed order. All local amplitudes have bounded derivatives of every
fixed order. The centered gamma cumulants are
\[
 \kappa_j(U)=(j-1)!\frac{t}{r^{j-1}}\quad(j\geq2),\qquad\kappa_1(U)=0.
\]
Consequently, Taylor expansion through degree $2J+1$, with its
$2J+2$ remainder, yields \eqref{joint:eq:all-orders} initially with
coefficients that are $e^{\delta\lambda}$ times polynomials in
$\lambda,t,t^{-1}$. The moment--cumulant formula shows that each
omitted contribution is $O(t^{J+1}m^{-J-1})$. The signed odd moments
are again used.

To identify these coefficients without retaining spurious negative
powers of $t$, introduce the auxiliary holomorphic family
\[
 \mathcal R(\xi)=\frac{r^{n+1}}{n!}
   \int_0^\infty
      \bigl(y+g(\xi e^{-y})\bigr)^n C(\xi e^{-y})^m e^{-ry}\,dy.
\]
At $\xi=1$ this is the original normalized integral. Choose a fixed
$a_0>0$ small enough that $g$ and $C$ are holomorphic on a
neighborhood of $|x|\leq a_0$. For $|\xi|\leq a_0$ the displayed
family is holomorphic; the exponents $n,m$ are integers, so zeros
of the factors introduce no branch choices.

The local expansion is uniform on that complex $\xi$-disk, with
$\xi\lambda$ replacing $\lambda$. Here is the necessary tail check,
since positivity cannot be used for complex $\xi$. For $y\geq1$,
\[
 \left|1+\frac{g(\xi e^{-y})}{y}\right|^n
       |C(\xi e^{-y})|^m
 \leq\exp\left(Kme^{-y}+Kne^{-y}/y\right).
\]
On $y\geq t/2$ this is bounded by $\exp(K'me^{-y})$, so the
previous density argument applies. On $Y_0\leq y\leq t/2$,
choose the fixed $Y_0$ sufficiently large: the last exponent is
then at most $rt/16$, whereas the gamma density has a deficit of
at least $rt/8$. On $0\leq y\leq Y_0$, bound the original
integrand's two analytic factors by $L^nK_0^m$, with fixed
$L,K_0$. After multiplication by $r^{n+1}/n!$, its logarithm is
$rt(1-\log t+\log L)+O(m+\log n)$, again negligible.
Thus the all-order remainders are uniform for $|\xi|\leq a_0$.
Cauchy's coefficient formula consequently permits extraction of any
fixed coefficient of $\xi^j$ from those finite asymptotic expansions.

For such a fixed $j$, take $n\geq j$, which holds eventually in
our regime. The binomial expansion is finite, and the elementary
integral
\[
 \binom nh\frac{r^{n+1}}{n!}
 \int_0^\infty y^{n-h}e^{-(r+j)y}\,dy
 =\frac1{h!}\left(\frac r{r+j}\right)^{n+1}(r+j)^h
\]
shows that the exact $\xi^j$ coefficient is
\begin{equation}\label{joint:eq:coefficient-integration}
 [\xi^j]\mathcal R(\xi)
   =\left(\frac r{r+j}\right)^{n+1}Q_j(m).
\end{equation}
The restriction $n\geq j$ ensures that every power of $g$ that can
contribute is present in the finite binomial expansion.
These operations concern a fixed Taylor coefficient; no divergent
infinite sum is integrated.

Now set $z=1/m$. The exact identities
\[
 z^jQ_j(1/z)=[w^j]K_j(z,w),\qquad
 e^{jt}\left(\frac r{r+j}\right)^{rt}=e^{\Theta_j(z,t)}
\]
show that the expansion of \eqref{joint:eq:coefficient-integration}
is exactly the $j$th coefficient in
\eqref{joint:eq:polynomial-prescription}, with $\xi\lambda$ in place
of $\lambda$. Compare successive powers of $m^{-1}$.
Their coefficient functions are Laurent polynomials in $t$;
a nonzero such function cannot be absorbed by an extra factor
$m^{-1}$ when $t=\log m+O(1)$. Thus uniqueness identifies every
fixed $\xi^j$ coefficient at every fixed asymptotic order.
The local coefficient functions are entire in $\xi\lambda$:
they are an exponential $e^{\delta\xi\lambda}$ times a polynomial
in $\xi\lambda$. Equality of all their Taylor coefficients identifies
them as entire functions, permitting evaluation at $\xi=1$ without
requiring analytic continuation of the full auxiliary integral to
that point. This proves the finite recipe and the cancellation of
every negative power of $t$.
It completes the all-order part of the theorem.

\paragraph{An independent derivation of $P_1$ and $P_2$.}
Let $E=e^{\delta\lambda}$, and write $D=\lambda\,d/d\lambda$.
The first two subleading offset-residue sums are
\begin{align*}
 T_1={}&h_2\lambda^2E+g_1(D+1)(\lambda E),\\
 T_2={}&h_3\lambda^3E+\frac{h_2^2}{2}\lambda^4E
   +g_2(D+1)(\lambda^2E)
   +g_1h_2(D+1)(\lambda^3E)
   +\frac{g_1^2}{2}(D+1)^2(\lambda^2E).
\end{align*}
Formula \eqref{joint:eq:polynomial-prescription} yields
\begin{align*}
 EP_1&=T_1+\frac t2D^2E,\\
 EP_2&=T_2+\frac t2D^2T_1
       +\left(-\frac t2D^2-\frac t3D^3+\frac{t^2}{8}D^4\right)E.
\end{align*}
These give exactly \eqref{joint:eq:P1}--\eqref{joint:eq:P2}; the
accompanying exact symbolic verifier checks agreement with the
independent gamma-moment derivative calculation.

\subsection{Numerical diagnostics and the scope of the advance}

The script \texttt{verify\_joint\_moments.py} evaluates the defining
integral after $x=e^{-y}$ and then centers and standardizes $y$ using
the gamma mean and variance. It does not evaluate a residue expansion.
It compares the leading transition function, one correction, and two
corrections at three offsets $s=-1,0,1$ and several increasing $m$.
The two integration-window choices are also compared.
These high-precision checks are diagnostics, not interval certificates;
the localization argument and Taylor remainders are the proof.

For a compact example, define
$A_J=e^{\delta\lambda}\sum_{\ell=0}^JP_\ell(t,\lambda)/m^\ell$.
The following entries use
$n=\operatorname{round}((m+1)\log m-1)$ and the exact $t,\lambda$
from the theorem. All displayed decimals are rounded numerical
approximations, including the errors.
\begin{center}
\begin{tabular}{rrrcc}
\hline
$m$ & $n$ & $\mathcal R_{n,m}$ & $|\mathcal R_{n,m}-A_1|/\mathcal R_{n,m}$
   & $|\mathcal R_{n,m}-A_2|/\mathcal R_{n,m}$\\
\hline
100 & 464 & 2.3971135731 & $1.1853\times10^{-3}$ & $1.1477\times10^{-4}$\\
300 & 1716 & 2.3595250117 & $2.3130\times10^{-4}$ & $8.0445\times10^{-6}$\\
1000 & 6914 & 2.3435725531 & $3.4368\times10^{-5}$ & $4.1120\times10^{-7}$\\
3000 & 24026 & 2.3382795613 & $5.5591\times10^{-6}$ & $2.5343\times10^{-8}$\\
\hline
\end{tabular}
\end{center}
The calculation uses 80 decimal digits. Across all 15 cases in the
JSON data, changing the integration window from 25 to 30 standard
deviations changed the result by at most $6.3\times10^{-63}$.
That comparison is a reproducibility check, not a rigorous bound
for the quadrature error or an additional proof of the theorem.

The regime in this theorem is asymmetric: $n\sim m\log m$, with
$n$ the larger exponent. Symmetry supplies the corresponding result
when the roles of $n$ and $m$ are exchanged. Section~\ref{gs:sec:main} treats the complementary proportional regime
$n/m\to c\in(0,\infty)$, proving a unique nondegenerate interior saddle.
Neither compact-parameter theorem by itself provides a
uniform sharp-truncation result when the reflected exponent grows
on the inverse-gamma least-term scale.

\paragraph{Editorial correction to the research status.}
The pinned fixed-$m$ statements remain valid. Their warnings against
substituting a growing $m$ are necessary, as
\eqref{joint:eq:limit} demonstrates quantitatively. After integration
of this contribution, the open problem should be narrowed from
``the simultaneous large-exponent problem'' to uniform bridges beyond compact
$m\exp[-(n+1)/(m+1)]$ and compact proportional ratios,
using also Theorem~\ref{gs:thm:proportional}. No contradiction to an existing theorem is
claimed.

\subsection{Further research questions}

\begin{enumerate}
\item Prove a uniform continuation of the transition when
$\lambda=me^{-t}$ also tends to zero or infinity. The compact-$\lambda$
proof identifies its expansion but does not give useful relative
bounds once the polynomials in $\lambda$ grow with $m$.
\item Replace the finite interval certificate in
Theorem~\ref{gs:thm:monotonicity} by a direct analytic inequality,
and determine the best global lower bound for the logarithmic
slope ratio. The proportional-saddle uniqueness question itself
is settled in the next section.
\item Build a uniform approximation joining the interior saddle to
$x\asymp1/m$. The present theorem supplies the boundary value that
such an approximation must reproduce.
\item Study the large-$\ell$ behavior of the explicitly computable
$P_\ell$. Does their growth lead to a useful optimal-truncation law
inside this transition, and how does it interact with the already
proved inverse-gamma late coefficients at fixed $m$?
\item Determine the maximal growth of $m$ for which the manuscript's
half-first-omitted-term law remains uniform. The local amplitude
$B(-\tau)^m$ changes the saddle when $m$ grows, so a fixed-amplitude
argument is insufficient.
\item Extend the exact polynomial recipe to other analytic endpoint
pairs $f(x)=-\log x+g(x)$ and $B(x)=b_1x+O(x^2)$. The leading
transition constant is suggested by $g_1+b_2/b_1$; a general theorem
should state explicit global endpoint hypotheses, not only a formal
local series.
\end{enumerate}
